\section{Introduction}\label{sec:intro}

Global optimization asks for a provably optimal solution of a nonconvex
problem, not a locally optimal or approximately optimal one. For quadratic
and polynomial objectives with continuous and integer variables this is hard
in the worst case, even under strong structural restrictions. Minimizing a
quadratic function with a single negative eigenvalue over a polytope is
NP-hard~\citep{pardalos1991-quadratic-programming-with-one}, and box-constrained quadratic
programming remains strongly NP-hard when its interaction graph has
treewidth two~\citep{pia2026-treewidth-and-the-complexity-of}. An exact
method must also handle degenerate configurations: ties between distant
solutions, flat optimal faces, and optimal values very close to a decision
threshold. The analysis below treats each of them explicitly. Spatial
branch-and-bound solvers handle many instances well, but they have no
polynomial guarantee.

Smoothed analysis~\citep{spielman2004-smoothed-analysis-of-algorithms-why} measures the cost of an algorithm
on a worst-case instance after a small random perturbation of its data. For
discrete optimization with a linear objective, polynomial smoothed
complexity is closely tied to pseudo-polynomial
solvability~\citep{beier2006-typical-properties-of-winners-and,roglin2007-smoothed-analysis-of-integer-programming}.
For continuous problems, a linear tilt removes degeneracy qualitatively: for
semialgebraic problems under a constraint qualification, generic tilts give
a unique minimizer with quadratic
growth~\citep{lee2017-generic-properties-for-semialgebraic-programs}. Such
statements give no probability bounds, and an exact algorithm in the bit
model needs more. It needs quantitative probability bounds, a finite rational law
that a Turing machine can sample, and an exact way to finish the draws that
remain degenerate, which under a finite law have positive probability.

This paper studies exact global optimization under one random linear
perturbation. A rational base instance with objective $F_0$ and feasible set
$X$ is given. A vector $\gamma$ is drawn once from a finite rational law that
is computed from the base instance before sampling, and the task is an exact
global minimizer of the sampled objective
\[
 F_\gamma(x)=F_0(x)+\gamma^{\mathsf T}x\qquad\text{on }X .
\]
We ask for algorithms that are correct on every draw, never resample, and
have expected bit complexity polynomial in the input length for fixed
structural parameters, with explicit dependence on numerical ratios between
curvature, widths, and the noise scale $\sigma$. \Cref{sec:model} makes this
model precise.

\subsection{Main results}\label{sec:intro:results}

The results follow four structural routes. They use different structure,
and partly different noise models. \Cref{tab:intro:results} lists the main
theorems.

\emph{Few negative curvature directions} (\cref{sec:qp}). For a rational
quadratic program on a bounded polytope whose Hessian has $k$ negative
eigenvalues, the nonconvexity lies in a $k$-dimensional factor. The
algorithm searches $k$ auxiliary coordinates that carry the negative
curvature and solves a convex quadratic program at each query. Under an
independent Gaussian-like finite law on every original coefficient, the
expected bit work is $f(n_z,k,1+\nu\diam(\mathcal P)/\sigma)\poly(I)$ with an
absolute exponent, also for bounded mixed-integer programs with $n_z$
integer variables; here $\nu$ is the largest negative curvature and
$\mathcal P$ the polyhedral relaxation (\cref{thm:qp:gauss}). Noise aligned
with the negative directions gives a similar fixed-parameter bound for
continuous problems (\cref{thm:qp:aligned}). Under independent uniform noise
on every coefficient the bound is polynomial for each fixed $k$ but contains
a power of the dimension that grows with $k$ (\cref{thm:qp:uniform}), and
\cref{thm:lim:ambient} shows that the counted quantities really depend on
the dimension. For $k\le2$ a different argument gives a bound linear in a
numerical ratio (\cref{thm:qp:two}). Convex piecewise-quadratic unaries with
rational breakpoints minus a concave term with $k$ supplied factor rows allow
arbitrarily many integer variables
(\cref{thm:qp:sep,thm:qp:sep-gauss}); on integer boxes with separable convex
polynomial terms, a lattice of attainable objective values gives exactness
without any fallback (\cref{thm:int:lowrank}).

\emph{Sparse interactions} (\cref{sec:sp,sec:con}). For a fixed-degree
polynomial on a mixed box whose monomials fit the bags of a supplied tree
decomposition with at most $p$ coordinates per bag, the expected bit work is
\[
 C_0^p\Bigl[4+\frac{(1+n/2)Lw_{\max}}{2\sigma}\Bigr]^p\poly_d(I)
\]
under independent uniform noise on all coordinates, where $L$ bounds the
coordinate second derivatives and $w_{\max}$ the coordinate widths
(\cref{thm:sp:main}). The integer dimension and the negative inertia are
unrestricted, and no growth or nondegeneracy assumption is made. The bound is
polynomial for each fixed $p$, but it is not fixed-parameter in $p$, and
\cref{thm:lim:width} shows that its factor $n^{\Omega(p)}$ is necessary for
the retention rule used. The method extends to equality constraints given by
explicit polynomial maps or monotone implicit equations, to affine state
recurrences with polynomial actuators, to products of simplices, and to
order constraints $x_i\le x_j$ with binary variables
(\cref{thm:con:graph,thm:con:actuator,thm:con:simplex,thm:con:order}). For
arbitrary sparse affine constraints no such result holds unless
$\mathrm{NP}\subseteq\mathrm{ZPP}$ (\cref{thm:lim:constraints}).

\emph{A small core with tractable recourse} (\cref{sec:rec,sec:int}).
Suppose that fixing $k$ continuous core coordinates leaves a residual
problem that an oracle solves exactly, or with certified error, on every
residual subbox: a convex quadratic or polynomial problem, a quadratic
program whose residual interaction graph is a forest, or a native-integer
problem with an exact oracle under bound tightening. Then the expected bit
work of the core search is at most $8^k[3+(1+k)L/(2\sigma)]^k$ times a
polynomial in the input length that includes the oracle cost, where $L$
bounds the curvature of the core coordinates only; integer recourse adds an
exact solve of a small polynomial problem of cost $c_d^k$ times a polynomial
(\cref{thm:rec:qp,thm:rec:poly,thm:int:native}). The number of residual
variables, their coupling, and their curvature enter only through the input
length. Under noise on all coordinates, residual convexity suffices. Under
noise on the core alone, our continuous point theorem uses a certified uniform
strong-convexity modulus of continuous recourse (\cref{thm:rec:core-only}),
or, for separable convex integer flows and totally unimodular recourse, a
certificate that describes all tied optimal flows at once
(\cref{thm:int:face,thm:int:flow-boundary,thm:int:tu}).
These are sufficient conditions. The companion results described below
also provide point oracles for residual-convex cubics and supplied
convexifiers; general arithmetic obstructions start at degree four
(\cref{sec:lim:points}).

\emph{Strong noise} (\cref{thm:int:strong-field}). Without any width or core
parameter, if the noise scale dominates the variation of every coordinate
derivative, weighted by the size of the coordinate's domain and the degree
of the interaction graph, then most coordinates are pinned at a bound by the
sign of their derivative, and the others form small random components of
the interaction graph. Each component is solved exactly by a fixed-degree
algebraic procedure (\cref{thm:int:solver}). Here the noise is large, not
small, and the result describes strongly random linear parts rather than
slight perturbations. The sufficient scale \eqref{eq:int:sf-regime} is
conservative: for nonquadratic continuous components it contains the large
effective solver base $c_d$, whereas the quadratic case uses the smaller
weight $a_i=3$.

\begin{table}[tbp]
\centering
\footnotesize
\caption{Main bounds, exact on every draw of one finite law. Here $k$ is
negative inertia, supplied factor row count, or core size; $p$ is bag size;
$b$ is grid bits or Gaussian accuracy;
$Q_{\rm ex}=[3+(1+k/2)L/(2\sigma)]^k$ and
$Q_{\rm ap}=[3+(1+k)L/(2\sigma)]^k$. FPT parameters include numerical
ratios. Usual outputs are listed; rare fallback outputs can be exponentially
long.}
\label{tab:intro:results}
\begin{tabularx}{\textwidth}{@{}>{\raggedright\arraybackslash}p{0.12\textwidth}
>{\raggedright\arraybackslash}p{0.19\textwidth}
>{\raggedright\arraybackslash}p{0.1\textwidth}
>{\raggedright\arraybackslash}p{0.12\textwidth}
>{\raggedright\arraybackslash}X
>{\raggedright\arraybackslash}p{0.14\textwidth}@{}}
\toprule
Result & Structure & Noise & $b$ & Expected bit work & Output\\
\midrule
\multicolumn{6}{@{}l}{\emph{Few negative directions}}\\
\Cref{thm:qp:gauss} & QP on a mixed polytope & Gaussian-like & $\poly(I)$ &
$f(n_z,k,1+\nu\diam(\mathcal P)/\sigma)\poly(I)$ & rational\\
\Cref{thm:qp:aligned} & QP on a polytope & aligned & $\poly(I)$ &
$C_0^k(1+\mathcal H_{\rm al})\poly(I)$, FPT in $k$ and factor ratios & rational\\
\Cref{thm:qp:uniform} & QP on a mixed polytope & uniform & $\poly(I)$ &
polynomial for fixed $k,n_z$ and numerical scales; power of $n$ grows with $k$ & rational\\
\Cref{thm:qp:two} & QP on a polytope, $k\le2$ & uniform or Gaussian-like &
$\poly(I)$ & $\poly(I)(1+\nu S/\sigma)$ & rational\\
\Cref{thm:qp:sep,thm:qp:sep-gauss} & convex piecewise-quadratic unaries,
rational knots, plus a concave factor, any $n_z$ & aligned, uniform, Gaussian-like & $\poly(I)$ & law-dependent
factors of the cited theorems; no factor for $n_z$ & rational\\
\Cref{thm:int:lowrank} & integer box, convex piecewise-polynomial unaries
plus a concave factor & aligned & $\poly(I)$ & $8^kH_{\rm rat}\poly(I)$, no
fallback & rational\\
\midrule
\multicolumn{6}{@{}l}{\emph{Sparse interactions}}\\
\Cref{thm:sp:main} & degree-$d$ polynomial on a mixed box, bag size $p$ &
uniform & $\poly_d(I)$ & $C_0^p\Theta^p\poly_d(I)$ with
$\Theta=4+(1+n/2)Lw_{\max}/(2\sigma)$ & implicit; rational for QP\\
\Cref{thm:con:graph,thm:con:actuator,thm:con:simplex,thm:con:order} &
graphs, actuator recurrences, simplices, orders with binaries & uniform &
$\poly_d(I)$ & polynomial for fixed $p$ and displayed numerical scales & charted or implicit\\
\midrule
\multicolumn{6}{@{}l}{\emph{Small core with tractable recourse}}\\
\Cref{thm:rec:qp} & quadratic, exact recourse on every subbox & uniform &
$\poly(I)$ & $8^kQ_{\rm ex}\poly(I)$ & rational\\
\Cref{thm:rec:poly} & certified recourse, e.g.\ convex residual & uniform &
$\poly_d(I)$ & $8^kQ_{\rm ap}\poly_d(I)$ & implicit\\
\Cref{thm:rec:core-only} & residual Hessian $\succeq\mu I$ & core only &
$\poly_d(I)$ & $8^kQ_{\rm ap}\poly_d(I)$ & implicit\\
\Cref{thm:int:native} & native-integer recourse, exact oracle & uniform &
$\poly_d(I)$ & $[8^kQ_{\rm ex}+c_d^k]\poly_d(I)$ & label $+$ algebraic, size
$c_d^k\poly_d(I)$\\
\Cref{thm:int:flow-interior}, \cref{cor:int:bilinear} & convex integer flows,
interior cores or bilinear coupling & core only & $\poly_d(I)$ &
$[8^kQ_{\rm ex}+c_d^k]\poly_d(I)$ & label $+$ algebraic\\
\Cref{thm:int:flow-boundary,thm:int:tu} & convex integer flows, totally
unimodular systems & core only & $f_d(k)\poly_d(I)$ &
$f_d(k)Q_{\rm ex}\poly_d(I)$ & label $+$ algebraic, size
$f_d(k)\poly_d(I)$\\
\midrule
\multicolumn{6}{@{}l}{\emph{Strong noise}}\\
\Cref{thm:int:strong-field} & degree-$d$ polynomial on a mixed box & strong
uniform & any & $\poly_d(I+b)$ times a subcritical branching factor &
components; rational for QP\\
\bottomrule
\end{tabularx}
\end{table}

In every route the output is an exact optimizer of the sampled instance. For
quadratic objectives on rational polyhedral and box domains it is rational;
an implicit graph lift can be irrational. For polynomial objectives the usual
output is a compact implicit descriptor: fixed integer labels and active
bounds, and a small box on which the objective is certified strongly convex,
so that its unique minimizer there is the global optimizer. The descriptor
supports evaluation of the optimizer and the optimal value to any requested
precision in time polynomial in the precision. Results that finish with a
small exact algebraic solve return all coordinates of the optimizer as
functions of one real algebraic number. The strong-noise route returns one
such representation per random component, with the optimal value as a sum
of component values. Bilinear flow coupling and interior optimal cores need
only polynomially many sampling bits (\cref{cor:int:bilinear,thm:int:flow-interior}).

\subsection{One mechanism}\label{sec:intro:mechanism}

The first three routes use one mechanism with four steps. Stating it once
explains what each theorem needs from its structure.

\emph{Curvature-corrected refinement preserves global optimizers.} Each route
has a search space, namely auxiliary factor coordinates, bag tuples, or core
coordinates, and a value function on it obtained by minimizing over the other
variables. Minimizing over a fixed feasible set preserves an upper bound on
coordinate curvature. Hence the minimum of the value function over a grid
cell is at least its least corner value minus a correction proportional to
the curvature bound and the squared mesh. Refining only the cells whose
corrected bound does not exceed the incumbent never discards a global
optimizer. This holds on every draw, including draws with ties.

\emph{Independent linear noise bounds the expected number of retained
states.} A retained cell has a corner whose value is within a fixed multiple
of the sum of the correction and certified evaluation error of the optimum.
Comparing such a corner with its two grid
neighbors in one coordinate confines that coordinate's noise coefficient to
an interval whose length is proportional to the curvature bound times the
mesh. The interval does not depend on the other coefficients of the search
space, so when these coefficients are independent the probabilities
multiply. Summing over a deterministic grid gives an expected count that
does not depend on the mesh. The count is taken over a grid fixed before
sampling, so no independence between the noise and the adaptively retained
cells is assumed. The direct product argument applies to sparse box bags,
cores, and aligned noise; constrained-domain variants use their proved
directional and transport bounds. Under uniform noise on all
coefficients of the low-rank route they are not, because the factor
coefficients depend on the residual noise; a conditional volume estimate
(\cref{lem:qp:volume}) replaces independence there and costs a power of
the dimension, and Gaussian-like noise restores independence in the
continuous proxy.

\emph{Exact closure finishes the sampled problem.} Refinement alone gives
approximations. An exact answer comes from a certificate proving that a
small exactly solvable problem contains a global optimizer: a region on
which a parametric convex quadratic program has one formula, a strongly
convex patch after fixing integer labels and active bounds, a gap that
excludes every competing integer label, or a certificate that fixes the face
of the core and the set of tied optimal flows. Each certificate is checked
exactly and is sound on every draw.

\emph{A base-selected finite law pays for an exact fallback on the same
draw.} The certificate passes by a refinement depth computed from base data
unless the draw lies in a rare exceptional set: small growth, a small
gradient at an active bound, a near-tie between integer labels, or noise
near one of finitely many fixed hyperplanes. The finite law is chosen fine
enough that this set has probability at most the reciprocal of the cost
factor of an exact fallback algorithm, which then runs on the same draw. The
cost factor is fixed first, then the thresholds and the depth, and the law
last, so that no precision depends on the sampled coefficients it is meant
to control.

The common invariant is a valid curvature-corrected bound together with a
noise comparison whose deterministic part excludes the tested coefficient.
Direct recourse obtains it by minimizing over a feasible set independent of
the searched coordinates. Graph charts and order transport replace that
literal premise with their proved curvature and feasible-comparison
arguments. The boundaries in \cref{sec:lim} delimit these mechanisms. Two
results finish differently: the integer lattice theorem
(\cref{thm:int:lowrank}) stops exactly once a certified gap falls below the
spacing of the attainable objective values, and the strong-noise theorem
(\cref{thm:int:strong-field}) solves its random components exactly on every
draw. Neither needs a rare fallback.

\subsection{What the guarantees mean}\label{sec:intro:meaning}

Several distinctions matter for reading the theorems.

\emph{Parameters.} A fixed-parameter bound has the form $f(\kappa)\poly(I)$
with an exponent independent of $\kappa$; the low-rank results under aligned
or Gaussian-like noise and the core results are of this form, with the
numerical ratio included in $\kappa$. The sparse results are polynomial for
each fixed width but are not fixed-parameter in the width. A bound that is
polynomial in a ratio such as $L/\sigma$ is polynomial in its magnitude.
For fixed structure, polynomially bounded ratios give polynomial work, but
$(1+L/\sigma)^{O(k)}\poly(I)$ with $L/\sigma=\poly(I)$ gives $I^{O(k)}$;
it does not give FPT in $k$ alone.

\emph{Every draw versus expectation.} Correctness holds on every draw,
including ties and draws with a continuum of optimizers. Only the work bound
is an expectation, over all draws; on rare draws the work can be exponential.

\emph{The law.} The finite law is constructed from the base instance; its
grid size or accuracy is chosen after a fallback cost factor and a stopping
depth are known. The theorems do not cover an arbitrary finite law, a coarse
grid chosen independently of the instance, or an arbitrary approximation of
a Gaussian. Most results use polynomially many random bits per coefficient,
and for them one resolution depending only on the input length suffices
(\cref{prop:model:uniform}); the nonlinear integer-flow results with
arbitrary core faces use a number of bits that depends on the core size.

\emph{Outputs.} Rational outputs are explicit. Implicit outputs are compact
descriptors whose global optimality is established by a separate proof
record; the descriptor has polynomial length, the expected length of the
record obeys the corresponding structural and numerical work bound, and
evaluation to $q$ bits costs time
polynomial in $I+q$. Algebraic outputs from a fallback can be exponentially
long on rare draws. No result promises expanded minimal polynomials of small
degree or exact comparisons of sums of unrelated algebraic numbers. No
route promises a canonical choice among tied optimizers on every draw.

\emph{The original objective.} All theorems optimize the sampled objective.
\Cref{prop:model:regret} turns a sampled solution into a certified interval
for the original optimum, of length at most the width of the realized
perturbation over $X$, $\max_X\gamma^{\mathsf T}x-\min_X\gamma^{\mathsf T}x$,
plus the evaluation error. Choosing the noise scale so that the largest
possible such width is at most $\varepsilon/2$, which for Gaussian-like laws
must account for their support radius (\cref{sec:model:original}), for routes valid at every positive scale, gives
on every draw an exactly feasible point whose original objective value is
certified to be within $\varepsilon$ of optimal. The expected work is
polynomial in $1/\varepsilon$ when the structure and the other numerical
parameters are fixed, and polynomial in $I$ and $1/\varepsilon$ only when
those parameters are polynomially bounded. On implicit graphs the
feasible point has rational free coordinates and an exact algebraic lift. We
do not claim that this improves on dedicated approximation algorithms, and
the strong-noise results give no such consequence. No result solves the unperturbed problem exactly,
which is consistent with the worst-case hardness cited above.

\subsection{Boundaries}\label{sec:intro:boundaries}

\Cref{sec:lim} delimits the method by explicit families. For the sparse
retention rule, connected quartic instances of treewidth $p-1$, with fixed
curvature, widths, and noise scale, retain more than $(5n/(6p))^{p/2}$ cells
on every draw, so the rule admits no bound $f(p)\poly(I)$
(\cref{thm:lim:width}); replacing the global rounding allowance by a bag-local
one deletes all optimizers (\cref{prop:lim:local}). Certified conditional
values give one sound way to remove the ambient dimension from this count
(\cref{prop:sp:conditional,prop:lim:conditional}). Under uniform ambient
noise, perfectly conditioned low-rank instances have expected counts growing
like $(N/k)^{k/4}$ (\cref{thm:lim:ambient}). Objective noise does not make
arbitrary sparse affine constraints of treewidth three tractable unless
$\mathrm{NP}\subseteq\mathrm{ZPP}$, because every exact sampled optimizer of
a suitable family decides \textsc{Subset Sum} (\cref{thm:lim:constraints}).
Under core-only noise with convex quartic
residual problems, a constant-accuracy residual point evaluator with
expected polynomial work would give Las Vegas algorithms for Square Root Sum
and $\PosSLP$. The Square Root Sum implication already holds at treewidth
two with bounded coefficients; the $\PosSLP$ construction has neither
restriction (\cref{thm:lim:posslp}). Value accuracy with polynomially
many bits does not localize residual points (\cref{prop:lim:value}). Single-flow
certificates fail at boundary core optima on an event of probability
$1/16$ (\cref{prop:lim:flow}). These are limitations of specified mechanisms
or conditional implications, not unconditional lower bounds, and several of
the instances are easy for other methods.

\subsection{Relation to prior work}\label{sec:intro:prior}

The ingredients of the mechanism are classical, and we use them as such:
semiconcavity and corner rounding on grids, nonserial dynamic programming on
tree decompositions~\citep{bertele1972-nonserial-dynamic-programming}, exact polynomial-time convex
quadratic programming~\citep{kozlov1980-the-polynomial-solvability-of-convex}, convex optimization
in the bit model~\citep{grotschel1988-geometric-algorithms-and-combinatorial-optimization}, exact convex
mixed-integer quadratic programming in fixed integer
dimension~\citep{pia2025-convex-quadratic-sets-and-the}, separable convex optimization over totally
unimodular systems~\citep{hochbaum1990-convex-separable-optimization-is-not}, quantifier elimination
with block-sensitive bounds~\citep{renegar1992-on-the-computational-complexity-and}
and other effective real-algebraic
methods~\citep{basu2006-algorithms-in-real-algebraic-geometry,basu2021-hausdorff-approximations-and-volume-of}, and isolation-type
anticoncentration for perturbed linear
objectives~\citep{beier2006-typical-properties-of-winners-and,roglin2007-smoothed-analysis-of-integer-programming}. Generic uniqueness
and growth under linear tilts of semialgebraic problems are known
qualitatively~\citep{lee2017-generic-properties-for-semialgebraic-programs}; \cref{sec:count} quantifies them for the
finite laws used here.

Among smoothed analyses of continuous global optimization,
\citet{kelner2007-on-the-hardness-and-smoothed} give expected polynomial time for minimizing
constant-rank quasi-concave functions over integral polytopes under a random
rotation of the objective subspace, by bounding the vertices of a projected
polytope. Our low-rank route perturbs the linear coefficients instead,
allows indefinite quadratic objectives that are not quasi-concave, and
returns exact rational optimizers. For problems with few negative
eigenvalues, deterministic algorithms give approximations whose work grows
polynomially in the inverse accuracy for fixed negative inertia~\citep{vavasis1992-approximation-algorithms-for-indefinite-quadratic}.
Rational Jacobi rotations also yield approximation algorithms for
mixed-integer quadratic programs with fixed integer dimension and negative
inertia~\citep{pia2026-rational-jacobi-rotations-and-the}.
For bounded treewidth, linear programming hierarchies give approximations of
polynomial optimization problems in time polynomial for fixed
width~\citep{bienstock2018-lp-formulations-for-polynomial-optimization}. The original-objective consequence of our
results (\cref{sec:intro:meaning}) is of the same type and is not claimed to
be better.

Recent deterministic exact results address other structured classes.
\citet{khajavirad2026-a-polynomial-time-solvable-class} solves sparse box
polynomial problems under sign, hidden-binary, width, and component-curvature
conditions; the quadratic result is in the rational bit model, whereas its
higher-degree result uses a unit-cost algebraic RAM. Exact integer and
mixed-integer quadratic algorithms of
\citet{ari2026-curvature-batching-for-integer-and} have
dimension- and magnitude-dependent bounds for integer data on polytopes.
\citet{ari2026-bounded-integer-quadratic-programming-through} treats bounded
pure integer QP, and \citet{wei2026-integer-quadratic-programming-in-fixed}
gives polynomial time for pure integer QP in each fixed integer dimension.
These premises and parameters differ from the arbitrary-dimensional
continuous residual problems and finite-noise expected bit work studied
here. Noise alone gives no general tractability principle:
\citet{soltanalian2026-random-parameter-noise-does-not} rules out uniform
expected polynomial bit time for exact ReLU verification under rounded,
clipped Gaussian parameter noise, assuming $\mathrm{NP}\nsubseteq\mathrm{BPP}$.
That result concerns network verification rather than the structured
optimization classes of our theorems.

Critical-region descriptions for parametric convex QP are established
machinery~\citep{tndel2003-an-algorithm-for-multi-parametric}; their existence
does not bound the number of regions polynomially. Likewise, classical
finite-set isolation and Pareto-count analyses concern discrete feasible
alternatives~\citep{mulmuley1987-matching-is-as-easy-as,beier2022-the-smoothed-number-of-pareto}.
Our continuous growth estimates and transfers to a fixed atomic law are
separate steps. Algebraic tube bounds control continuous volume; the finite
law requires the additional grid-transfer argument.

The results of this paper are specific certificates, bounds, and
compositions built from these ingredients. For few negative directions they are
closure by regions of a parametric convex program, with a fixed-hyperplane
exceptional set, and the resulting fixed-parameter expected bounds under
aligned and Gaussian-like noise and the dimension-dependent bound under
uniform noise (\cref{sec:qp}). For sparse interactions they are exact
closure on certified strongly convex patches after a global-error bag
search, with fixed-width expected bounds, and the extensions to graph,
recurrence, simplex, and order constraints (\cref{sec:sp,sec:con}). For small
cores they are the excluded-region certificate for continuous recourse, the
competing-label certificate for integer recourse, and the optimal-face
certificate for core-only integer flows and totally unimodular systems
(\cref{sec:rec,sec:int}). Further contributions are the strong-noise
component composition (\cref{sec:int}) and the boundary results of
\cref{sec:lim}. Exactness on every draw of a base-chosen finite law is
carried through all of these results, but we do not regard the finite-law
model itself, isolation-type anticoncentration, representations of
algebraic points by one root, or connected-set counting as new.

Three unpublished companion manuscripts have specific overlaps.
\emph{Exact Arithmetic in Polynomial Optimization: Values, Optimizers, and
Certificates}~\citep{companion-exact-arithmetic} supplies all-precision
value and selected-core oracles for core-noise problems with convex
residuals, including coupled-polytope versions and exact rational QP
reconstruction. It also supplies full selected-point oracles for
residual-convex cubics without a residual modulus or convexifier, and under
a supplied convexifier when the transformed objective is cubic on the
feasible polytope or is a fixed-degree polynomial convex on all of $\R^n$.
Its cubic selector is the lexicographically least optimal core followed by
the least-norm residual optimizer in that fiber. One base-selected law and
one per-draw work factor $\Omega(\gamma)$ apply to all precision queries
$q$ for this fixed point. These law and query properties already appear in
the companion. Its outputs are Cauchy names for the selected point and value;
our ordinary outputs are finite patch descriptors with a separate global
proof record, and need not select the same canonical point. This is a
representation distinction. For uniformly strongly convex residuals, the
direct recourse count also preserves the original $L/\sigma$, whereas a
derived Schur convexifier can require scale $M+M^2/\mu$
(\cref{sec:lim:rank}).
The finite-law transfer, section bounds, projected-growth argument, and
lexicographic fallback overlap with our foundations. Its quartic point
reductions, including their core-only-noise consequence, and regularization
example are reproduced with proofs in
\cref{thm:lim:posslp,prop:lim:value}. Thus quartic point obstructions do not
exclude its positive cubic result.

\emph{Decomposition-aware global optimization: certified coordinate grids,
conditional recourse, and structural limits}~\citep{companion-decomposition-aware}
develops deterministic curvature-corrected grids, conditional-recourse
filtering, boundary output, and the star obstruction. For its box quadratic
class, let $L_i=\max\{0,\sup_{\widehat X}\partial_{ii}F\}$ and
$\norm d_L^2=\sum_iL_id_i^2$. Weighted growth
$F(x)-\min F\ge\gamma_*\norm{x-x^*}_L^2$ gives
$\bar\kappa=\max\{1,1/\gamma_*\}$; point growth with modulus $g$ gives
$\kappa=\max\{1,\max_iL_i/g\}$ and $\bar\kappa\le\kappa$. The companion
certifies $2^{-q}$ accuracy in $f(p,\bar\kappa)(I+q+1)^5$ bit work, with
\[
 f(p,t)=c_0(c_1p\sqrt t)^p\,t(1+\log_2t)^2,
\]
and gives exact output in $f_1(p,\kappa)(I+1)^{C_1}$ under point growth at
a unique optimizer. Graded grids use
$O(\sqrt{\bar\kappa}\log(n+2))$ nodes per coordinate; its sharp example
for the filtered common uniform-grid method needs
$\sqrt{(n-2)\kappa}+1$ in the stated stages. Under rETH its lower bound
excludes a sublinear width exponent in the restricted product-time model
$f(p)(\kappa I)^{o(p)}$ for unique-minimizer integer box QP with a supplied
decomposition. It is a
deterministic approximation boundary, not a smoothed lower bound or a
matching $p/2$ exponent result. The limits of integrating the
growth-dependent work under noise are discussed in \cref{sec:sp:width}.
Our sparse theorem bounds expected counts without a supplied growth modulus
and closes exactly on every draw.

\emph{Sparse indicator quadratics: exact complexity and smoothed separator
messages}~\citep{companion-sparse-indicator} treats positive-definite
indicator quadratics and perturbs only binary indicator penalties by
independent uniform-grid draws, with at least $2n$ labels per marginal. It
returns exact answers and message dictionaries on a stated bounded boundary
box on every draw, with fixed-width expected work polynomial in input length
and the displayed numerical magnitudes. Removing those numerical factors
would give $\mathrm{NP}\subseteq\mathrm{ZPP}$, even at its fixed matrix
family and noise scale. These finite-noise, width, and dictionary ideas
overlap. Our sparse route perturbs every original
linear coefficient and treats continuous and native integer coordinates;
its affine-feasibility boundary does not contradict that indicator model.
The proofs needed here are included in this paper. The bounds are
mathematical guarantees; practical performance is unestablished.

\subsection{Organization}\label{sec:intro:organization}

\Cref{sec:model} fixes the model and output contracts. \Cref{sec:count}
develops the shared counting, tail, and fallback tools. \Cref{sec:qp} treats
few negative directions, \cref{sec:sp,sec:con} sparse interactions,
\cref{sec:rec,sec:int} small cores with continuous and integer recourse and
the strong-noise route. \Cref{sec:lim} proves the boundary results, and
\cref{sec:disc} discusses uses and open questions. Appendices contain the
longer proofs.
